\documentclass[10pt,a4paper]{amsart}
\usepackage[margin=19mm]{geometry}
\usepackage{amsmath,amssymb,mathtools}
\usepackage[scr=rsfs,cal=euler]{mathalfa}
\usepackage{xfrac}
\usepackage[unicode,hidelinks,bookmarks=false]{hyperref}
\newcommand{\ie}{i.e.\ }
\newcommand{\cf}{cf.\ }
\newcommand{\resp}{respectively\ }
\newcommand{\Subsection}{\subsection}
\providecommand{\nobreakdash}{\nobreak\hbox{-}}
\setlength{\emergencystretch}{2em}
\multlinegap0pt
\usepackage{bm}
\usepackage{bbm}
\usepackage{dsfont}
\usepackage{xspace}
\usepackage{cancel}
\usepackage{mathtools}
\usepackage{amsthm}
\makeatletter
\@ifundefined{theorem}{\newtheorem{theorem}{Theorem}}{}
\@ifundefined{definition}{\newtheorem{definition}[theorem]{Definition}}{}
\@ifundefined{proposition}{\newtheorem{proposition}[theorem]{Proposition}}{}
\makeatother
\newcommand{\RR}{\mathbb{R}}
\title[A billiard table close to an ellipse is deformationally spec]{A billiard table close to an ellipse is deformationally spectrally rigid among dihedrally symmetric domains}
\author{Corentin Fierobe, Vadim Kaloshin, Alfonso Sorrentino}
\date{Source: arXiv:2511.20062v3 (CC BY 4.0)}
\begin{document}
\maketitle

\noindent\emph{Source and licence.} A billiard table close to an ellipse is deformationally spectrally rigid among dihedrally symmetric domains. Version arXiv:2511.20062v3 (CC BY 4.0), distributed under the Creative Commons Attribution 4.0 International licence, as recorded in the arXiv licence metadata for this exact version and shown on the abstract page https://arxiv.org/abs/2511.20062v3. Full rights notice: licenses/arxiv-2511.20062-CC-BY-4.0.txt.\par\smallskip

\noindent\textit{Editorial scope.} Counted unit: the Proposition whose source label is \texttt{proposition:first\_var\_beta} in the article (source0.tex lines 662--669, source label \texttt{proposition:first\_var\_beta}), delivered together with the article's own complete original proof of it (source0.tex lines 672--741, one \texttt{proof} environment of 70 lines). No mathematics is written, repaired or re-derived: statement and proof are the article's own text line for line. Required context retained uncounted: equation:beta\_inv\_curve (lines 587--590); defntau (lines 606--612). Every definition, display and label of those lines is retained unchanged. Omitted from this package and not redistributed: 1--586, 591--605, 613--661, 670--671, 742--1910. Nothing inside a retained excerpt is omitted or altered: every retained body is delivered exactly as the article's lines are written, so no omission receipt of any kind accompanies this package. Citations: the retained excerpts carry no citation command at all, so no bibliography entry of the article is transcribed and no reference is redistributed. Reference adaptation: none was needed and none was made; every reference inside the delivered unit resolves inside it, so no unresolved reference or citation is produced. Printed slips kept literal: no printed word, symbol, hypothesis or number of the counted statement or of its proof was altered. Layout and class: the article's own class is \texttt{amsart} with packages including inputenc, fontenc, graphicx, amsmath, amssymb, amsthm, mathrsfs, hyperref, pgf, tikz, pgfplots, array, color, mathtools; the wrapper is the lane's pinned \texttt{amsart} editorial wrapper at 10pt on A4, which restates the theorem-like environments the retained text needs and adds the article's own macro definitions that the retained text needs (\texttt{\textbackslash RR}), with extra wrapper packages (bm, bbm, dsfont, xspace, cancel, mathtools). The delivered numbering is the unit's own. Assets: no retained span contains a figure environment, a graphics-inclusion command, a PDF-page inclusion, a table or a TikZ picture; no image file is copied into this package, the package declares no asset dependency and no image file is redistributed with it. Licence: version arXiv:2511.20062v3 of this article is distributed under the Creative Commons Attribution 4.0 International licence, as recorded in the arXiv licence metadata for this exact version and shown on the abstract page https://arxiv.org/abs/2511.20062v3. A full rights notice is delivered at licenses/arxiv-2511.20062-CC-BY-4.0.txt. Characters: every retained excerpt is pure ASCII, so the delivered engine, XeLaTeX with its default Latin Modern text font, reports no missing character.
\begin{equation}
\label{equation:beta_inv_curve}
\beta(\omega) = -\int_0^1|\gamma(s(\theta+\omega))-\gamma(s(\theta))|d\theta.
\end{equation}

\begin{definition}
The \textit{deformation map} associated to $(\Omega_{\tau})_{\tau\in I}$ is the family of maps $(n_{\tau})_{\tau\in I}$ defined for any $\tau\in I$ and $s\in\RR$ by
\begin{equation} \label{defntau}
n_{\tau}(s) = \langle\partial_{\tau}\gamma_{\tau}(s)\,,\,N_{\gamma_{\tau}}(s)\rangle,
\end{equation}
where $\langle\,,\,\rangle$ denotes the canonical scalar product on $\RR^2$ and $N_{\gamma_{\tau}}(s)$ is the outgoing unit normal vector to $\partial\Omega_{\tau}$ at the point $\gamma_{\tau}(s)$. 
\end{definition}

\begin{proposition}
\label{proposition:first_var_beta}
Assume that each domain $\Omega_{\tau}$ has a rotational invariant curve parametrized by $\Gamma_{\tau}$ of irrational rotation number $\omega$ such that the map $(\tau,\theta)\mapsto\Gamma_{\tau}(\theta)$ is $\mathscr C^1$-smooth and conjugates the billiard map in $\Omega_{\tau}$ to a rotation. Then the map $\tau\mapsto\beta_{\tau}(\omega)$ is differentiable and its $\tau$-derivative is given by
\begin{equation}
\label{equation:beta_form1}
\partial_{\tau}\beta_{\tau}(\omega) = 2\int_0^{1}n_{\tau}(s_{\tau}(\theta))\sin\varphi_{\tau}(\theta)d\theta.
\end{equation}
\end{proposition}

\begin{proof}%[Proof of Proposition \ref{proposition:first_var_beta}]
According to \eqref{equation:beta_inv_curve}, $\beta_{\tau}(\omega)$ admits the following expression:
\[
\beta_{\tau}(\omega) = -\int_0^1|\gamma_{\tau}(s_{\tau,\omega}(\theta+\omega))-\gamma_{\tau}(s_{\tau,\omega}(\theta))|d\theta,
\]
where $s_{\tau,\omega}(\theta)$ is the $s$-projection of $\Gamma_{\tau}(\theta)$, and $\gamma_{\tau}$ is the corresponding parametrization of $\partial\Omega_{\tau}$ by $s$.

To ease the computations, we drop the index $\omega$ in $s_{\tau,\omega}$. Let us consider the familiy of generating maps $L_{\tau}$, namely the maps defined for any $\tau$ by
\[
L_{\tau}(s_0,s_1) := -|\gamma_{\tau}(s_1)-\gamma_{\tau}(s_0)|,\qquad s_0,s_1\in\RR.
\]
Then
\begin{equation}
\label{equation:expression_beta_proof}
\beta_{\tau}(\omega) = -\int_0^{1}L_{\tau}(\gamma_{\tau}(s_{\tau}(\theta)),\gamma_{\tau}(s_{\tau}(\theta+\omega)))d\theta
\end{equation}

Differentiating \eqref{equation:expression_beta_proof} with respect to $\tau$ we obtain the formula
\[
-\partial_{\tau}\beta_{\tau}(\omega) = A+B+C
\]
where
\[
A := \int_0^{1} \partial_{\tau} L_{\tau}(\gamma_{\tau}(s_{\tau}(\theta)),\gamma_{\tau}(s_{\tau}(\theta+\omega)))d\theta,
\]
\[
B := \int_0^{1}
\partial_{\tau}\gamma_{\tau}(s_{\tau}(\theta)) \partial_{s_0}L_{\tau}(\gamma_{\tau}(s_{\tau}(\theta)),\gamma_{\tau}(s_{\tau}(\theta+\omega)))d\theta,
\]
and 
\[
C := \int_0^{1}
\partial_{\tau}\gamma_{\tau}(s_{\tau}(\theta+\omega)) \partial_{s_1}L_{\tau}(\gamma_{\tau}(s_{\tau}(\theta)),\gamma_{\tau}(s_{\tau}(\theta+\omega)))d\theta.
\]
Applying the change of coordinates $\theta'=\theta+\omega$ in the integral defining $C$, we obtain
\begin{eqnarray*}
C+B &=& \int_0^{1}
\partial_{\tau}\gamma_{\tau}(s_{\tau}(\theta))
\left( 
\partial_{s_1}L_{\tau}(\gamma_{\tau}(s_{\tau}(\theta-\omega)),\gamma_{\tau}(s_{\tau}(\theta)))\right. \\
&& \left. \qquad + \;
\partial_{s_0}L_{\tau}(\gamma_{\tau}(s_{\tau}(\theta)),\gamma_{\tau}(s_{\tau}(\theta+\omega)))
\right)
d\theta. 
\end{eqnarray*}
Since $L_\tau$ is a generating function for the billiard dynamics, then
\[
\partial_{s_1}L_{\tau}(\gamma_{\tau}(s_{\tau}(\theta-\omega)),\gamma_{\tau}(s_{\tau}(\theta)))
+
\partial_{s_0}L_{\tau}(\gamma_{\tau}(s_{\tau}(\theta)),\gamma_{\tau}(s_{\tau}(\theta+\omega))) \equiv 0.
\]
Hence, $B+C=0$. 

Now, if we set $s_0 = s_{\tau}(\theta)$, $s_1 = s_{\tau}(\theta+\omega)$, then
\[
\partial_{\tau} L_{\tau}(s_0,s_1) = 
\langle \partial_{\tau}\gamma_{\tau}(s_0), u^{\tau}(\theta) \rangle
-\langle \partial_{\tau}\gamma_{\tau}(s_1), u^{\tau}(\theta) \rangle,
\]
where $u^{\tau}(\theta)$ is the unit vector joining $\gamma_{\tau}(s_0)$ to $\gamma_{\tau}(s_1)$. Applying again the change of coordinates $\theta'=\theta+\omega$ while integrating the second term of $\partial_{\tau} L_{\tau}(s_{\tau}(\theta),s_{\tau}(\theta+\omega))$, we obtain
\[
\partial_{\tau}\beta_{\tau}(\omega) = \int_{0}^{1} 
\langle \partial_{\tau}\gamma_{\tau}(s_{\tau}(\theta)), u^{\tau}(\theta)-u^{\tau}(\theta-\omega) \rangle d\theta.
\]
Because of the billiard reflexion law off the boundary $\partial\Omega_{\tau}$ at the point $\gamma_{\tau}(s_{\tau}(\theta))$, 
\[
u^{\tau}(\theta)-u^{\tau}(\theta-\omega) = -2\sin\varphi_{\tau}(\theta)N_{\gamma_{\tau}}(s_{\tau}(\theta))
\]
and the claim then follows by using the definition of $n_\tau$ in \eqref{defntau}.
\end{proof}

\end{document}
